\documentclass{beamer}
% \documentclass[handout]{beamer}
% \documentclass[draft]{beamer}

% \usepackage{beamerthemebars}
% \usepackage{beamerthemeclassic}
% \usepackage{beamerthemelined}
\usepackage{beamerthemeshadow}
% \usepackage{beamerthemesidebar}
% \usepackage{beamerthemesplit}







\title{Computational investigation of intense short-wavelength laser interaction with rare gas clusters}
\subtitle{Ph.D. defence}
\author{Nicolas Bigaouette}
\date[December 2013]{\today}
\institute{Computational Nanophotonics Group
\newline
Physics Department -- University of Ottawa}
%\titlegraphic{\includegraphics[width=0.15\textwidth]{figures/UofO}}


\RequirePackage[utf8]{inputenc}
\usepackage[english]{babel}
\usepackage{dejavu}
\usefonttheme{serif}
\RequirePackage[T1]{fontenc}
%\usepackage{libertine}

\usepackage{color}

\usepackage{bm}				% Needed for bold Greek letters


%%\usepackage{textpos}		% http://www.latex-community.org/forum/viewtopic.php?f=47\&t=9387
%\usepackage[showboxes,absolute]{textpos} % http://handyfloss.net/2007.05/latex-the-textpos-package/
%%\usepackage{textpos} % http://handyfloss.net/2007.05/latex-the-textpos-package/

% https://sites.google.com/site/geraldinedelmondo/latex-1/trucs-et-astuces/placereltsbeamer
\usepackage[absolute,overlay]{textpos}
\setlength{\TPHorizModule}{\paperwidth}
\setlength{\TPVertModule}{\paperheight}

\newcommand{\framenumber}{
% Place text @ page relative
%					(x, y)
\begin{textblock}{0}(0.01,0.967)
\begin{scriptsize}
{\color{gray}\insertframenumber/\inserttotalframenumber}
\end{scriptsize}
\end{textblock}
}


%\usepackage{algorithmic}
%\usepackage{animate}
% \usepackage{movie15}

\newcommand{\pa}[1]{\left( #1 \right)}
\newcommand{\cro}[1]{\left[ #1 \right]}
\newcommand{\abs}[1]{\left| #1 \right|}
\newcommand{\ve}[1]{\textbf{#1}}
\newcommand{\vecr}{\ve{r}}
\newcommand{\vecnabla}{\bm{\nabla}}
\newcommand{\laplacien}[1]{\vecnabla^2 #1}           % Laplacien
\newcommand{\laplacient}[1]{\vecnabla_{\bot}^{2} #1} % Laplacien transverse
\newcommand{\gradient}[1]{\vecnabla #1}              % Gradient
\newcommand{\grad}[1]{\vecnabla #1}                  % Gradient
\newcommand{\divergence}[1]{\vecnabla \cdot #1}      % Divergence
\newcommand{\rotationnel}[1]{\vecnabla \times #1}    % Rotationnel
\newcommand{\vE}{\ve{E}}
\newcommand{\vB}{\ve{B}}
\newcommand{\vD}{\ve{D}}
\newcommand{\vH}{\ve{H}}
\newcommand{\vJ}{\ve{J}}
\newcommand{\del}{\partial}                                     % del
\newcommand{\delt}[1]{\frac{\del #1}{\del t}}                   % del ? / delt

\setbeamercovered{transparent}

\begin{document}

% ###############################################################
\begin{frame}{}
\titlepage
%\maketitle
% Place image @ page relative
%					(x, y)
\begin{textblock}{0}(0.04,0.6)
\includegraphics[width=0.22\paperwidth]{figures/UofO}
\end{textblock}
\end{frame}

\setcounter{framenumber}{0}

% ###############################################################
\begin{frame}\framenumber
\tableofcontents
\end{frame}


% ###############################################################
\section{Laser-Cluster interaction}

% ---------------------------------------------------------------
\subsection{Clusters of atoms and strong laser fields}
% ***************************************************************
\begin{frame}{}\framenumber
\end{frame}

% ---------------------------------------------------------------
\subsection{Mechanisms}
% ***************************************************************
\begin{frame}{}\framenumber
\end{frame}

% ...............................................................
\subsubsection{IR}
% ***************************************************************
\begin{frame}{}\framenumber
\end{frame}

% ...............................................................
\subsubsection{VUV \& XUV}
% ***************************************************************
\begin{frame}{}\framenumber
\end{frame}

% ...............................................................
\subsubsection{New processes}
% ***************************************************************
\begin{frame}{}\framenumber
\end{frame}


% ###############################################################
\section{Methodology and Tools}

% ---------------------------------------------------------------
\subsection{MD}
% ***************************************************************
\begin{frame}{}\framenumber
\end{frame}

% ---------------------------------------------------------------
\subsection{QFDTD}
% ***************************************************************
\begin{frame}{}\framenumber
\end{frame}


% ###############################################################
\section{Papers}

% ---------------------------------------------------------------
%\subsection{1. Augmented collisional ionization via excited states in XUV cluster interactions}
\subsection{1. ACI, XUV, Xenon: Introduction}
% ***************************************************************
\begin{frame}{}\framenumber
\end{frame}

% ---------------------------------------------------------------
%\subsection{2. Clusters in intense XUV pulses: Effects of cluster size on expansion dynamics and ionization}
\subsection{2. ACI, XUV, Argon: Cluster size}
% ***************************************************************
\begin{frame}{}\framenumber
\end{frame}

% ---------------------------------------------------------------
%\subsection{3. Effect of Augmented Collisional Ionization and potential depth in the VUV regime; a theoretical study}
\subsection{3. ACI, VUV, Xenon: 100nm experiment}
% ***************************************************************
\begin{frame}{}\framenumber
\end{frame}

% ---------------------------------------------------------------
%\subsection{4. Recombination effects in soft-x-ray cluster interactions at the xenon giant resonance}
\subsection{4. Recombination effects, Xenon}
% ***************************************************************
\begin{frame}{}\framenumber
\end{frame}

% ---------------------------------------------------------------
%\subsection{5. Nonlinear grid mapping applied to an FDTD-based, multi-center 3D Schr\"odinger equation solver}
\subsection{5. QFDTD}
% ***************************************************************
\begin{frame}{}\framenumber
\end{frame}


% ###############################################################
\section{Conclusion}

% ***************************************************************
\begin{frame}{}\framenumber
\end{frame}



\end{document}


